% LLM-assisted: Codex retained the numerical protocol and updated file
% locations for the submitted code inventory on 6 November 2026.
\section{Details for reproducing Methods}
\label{app:methods-reproduction}
This record gives the numerical ordering, transforms, random draws and
scoring calls used in the saved experiments. Current setup and execution
commands are in \nolinkurl{supplementary/REPRODUCIBILITY.md}.

\subsection{Required files and dataset release}
The submitted workflow uses \nolinkurl{Code/}, \nolinkurl{requirements-lock.txt},
\nolinkurl{data/polynomial_baselines.json} and the four experiment folders
under \nolinkurl{results/}. Download the raw AL-PUB archive separately.
Processed image caches can be rebuilt from that archive.

The raw file is \nolinkurl{data/alpub_v2/ALPUB_v2.zip}, downloaded on
27 October 2026 from \url{https://www.kaggle.com/datasets/miswindall/al-pub-v2}.
The dataset title is AL\_PUB\_v2. Its Kaggle release number is 1, dated
5 February 2024. The unchanged archive has 326979938 bytes and SHA-256
\begin{quote}\ttfamily\small
\nolinkurl{1d18415ed0875e10459c40e04e2a06c8c8c2036cbc78fbc2105f2da30c95da8a}
\end{quote}
The download endpoint is
\url{https://www.kaggle.com/api/v1/datasets/download/miswindall/al-pub-v2}.
Place the downloaded file at the path above and check its hash before running
an audit. On macOS, use:
\begin{verbatim}
shasum -a 256 data/alpub_v2/ALPUB_v2.zip
\end{verbatim}
The archive contains 205797 JPEGs. The audit excludes 21 blank crops and
keeps 205776 images. The excluded archive paths are recorded in
\nolinkurl{results/alpub_v2/duplicates.csv}, with the release counts and
split provenance in \nolinkurl{results/alpub_v2/audit.json}. The downloaded
release metadata is kept in \nolinkurl{data/alpub_v2/kaggle_metadata.json}.
Runge data are generated by the notebook and need no external download.


\subsection{Execution}
Follow \nolinkurl{supplementary/REPRODUCIBILITY.md} for environment setup,
synthetic checks and a full rerun. Regression runs through
\nolinkurl{Code/Project2_working.ipynb}. Classification uses the prepare,
search, comparisons and final stages of
\nolinkurl{Code/part_e_classification.py}. The notebook's final cells display
saved classification results, so classification runs first in a full rerun.

\subsection{Data ordering and numerical transforms}
Runge inputs used NumPy's \nolinkurl{linspace(-1,1,1000)} with endpoints
included. A fresh \nolinkurl{default_rng(42)} supplied Gaussian noise
through \nolinkurl{Generator.normal}, with mean zero and standard
deviation 0.1. The original indices $0,\ldots,999$ were passed to
\nolinkurl{train_test_split} in input order. We kept the returned
development order when applying \nolinkurl{KFold} to its positions, then
mapped fold positions to original row numbers. The split and fold
settings are stated in Methods. Feature statistics used
\nolinkurl{numpy.mean} and \nolinkurl{numpy.std} along columns, with
\nolinkurl{ddof=0}. Target centring used the training target mean.

For AL-PUB, JPEG archive paths were sorted lexicographically. Development
rows kept that order after filtering out excluded and test images.
Publication filenames were parsed for the POxy item number. Its group
string was \nolinkurl{POxy.n} followed by the integer item number,
without a volume number. Distinct group strings were sorted
lexicographically. NumPy's \nolinkurl{Generator.shuffle}, using
\nolinkurl{default_rng(2026)}, shuffled the 1419 group strings. We assigned
the first $\lceil0.2n_{\rm groups}\rceil=284$ groups to testing, rounding
20\% of the groups upwards. This procedure was first used in
\nolinkurl{Code/alpub_v2_feasibility.py}. The classification experiments reused the saved group assignments. Group separation was checked by
publication item number. Physical manuscript and writer identities were
unverified. The sample indices returned by \nolinkurl{StratifiedShuffleSplit} were sorted into
development-row order before making grouped folds. The fixed label order
is saved in \nolinkurl{class_mapping.json}.

Pillow converted each image to \nolinkurl{RGBA}, composited it onto opaque
white with \nolinkurl{alpha_composite}, then converted it to mode
\nolinkurl{L}. \nolinkurl{ImageOps.autocontrast} used its defaults.
\nolinkurl{thumbnail((26,26), Image.Resampling.LANCZOS)} kept the aspect
ratio. On the white $28\times28$ canvas, the left and top offsets were
$\lfloor(28-w)/2\rfloor$ and $\lfloor(28-h)/2\rfloor$, where $w,h$ are
the thumbnail dimensions. NumPy's default row-major \nolinkurl{ravel}
flattened the result. Cached uint8 pixels were converted to float64,
divided by 255, then subtracted from one.

\subsection{Random draws and scoring calls}
Initialisation and shuffle generators were separate fresh NumPy
\nolinkurl{default_rng} instances for each fit, with that experiment's
stated seeds. Weights were drawn layer by layer with
\nolinkurl{Generator.normal}. The shuffle generator's state continued
across epochs as \nolinkurl{Generator.permutation} supplied row orders.

The regression regularisation search used
\nolinkurl{ParameterSampler(n_iter=16, random_state=42)} with hidden-layer
list $[25]$, $[25,25]$, $[50,50]$ in that order. Keys were
\nolinkurl{hidden}, \nolinkurl{gamma} and \nolinkurl{lmbd}.
SciPy's \nolinkurl{loguniform(1e-4,1e-2)} supplied learning rates and
\nolinkurl{loguniform(1e-7,1e-2)} supplied positive regularisation parameters. The same
draws were used for both penalties. The unregularised family used the
same architecture and learning-rate draws with zero penalty.

Classification used \nolinkurl{ParameterSampler} with 12 draws and
\nolinkurl{random_state=2026}. The categorical lists were depth
$[1,2,3]$, width $[32,64,128]$, activation
\nolinkurl{['sigmoid','relu','leaky_relu']} and penalty
\nolinkurl{['none','l1','l2']}. The dictionary keys were
\nolinkurl{depth}, \nolinkurl{width}, \nolinkurl{activation},
\nolinkurl{gamma}, \nolinkurl{penalty} and \nolinkurl{positive_lmbd}.
Learning rates and regularisation parameters used the SciPy distributions stated in Appendix~\ref{app:experiment-settings}.
A positive strength was drawn for every hyperparameter configuration, then replaced by zero
when the penalty was \nolinkurl{none}.

Classification predictions used NumPy's row-wise \nolinkurl{argmax}.
Accuracy used scikit-learn's \nolinkurl{accuracy_score}. Macro-F1 used
\nolinkurl{f1_score} with labels $0,\ldots,23$ explicitly supplied,
\nolinkurl{average='macro'} and \nolinkurl{zero_division=0}.

\subsection{Scikit-learn FFNN}
The matched scikit-learn model used \nolinkurl{MLPClassifier} with
\nolinkurl{solver='adam'}, \nolinkurl{alpha=0},
\nolinkurl{batch_size=128}, \nolinkurl{shuffle=True} and
\nolinkurl{early_stopping=False}. Picked hidden widths and learning rates set
\nolinkurl{hidden_layer_sizes} and \nolinkurl{learning_rate_init}.
Sigmoid mapped to \nolinkurl{activation='logistic'}, and ReLU to
\nolinkurl{activation='relu'}. Each fit received a fresh
\nolinkurl{RandomState(2026)}. Each \nolinkurl{partial_fit} call supplied
all 24 classes and trained for one epoch. Other settings used the
defaults of the saved scikit-learn version.

\subsection{Software and saved experiment records}
The environment used Python 3.12.7 on macOS ARM64. The array calculations used NumPy 2.5.3 \parencite{harris2020numpy}.
Stable functions and sampled distributions used SciPy 1.18.1
\parencite{virtanen2020scipy}. Scikit-learn was 1.9.1, pandas was
2.3.3, Pillow was 12.3.0, and JAX/jaxlib were 0.11.2.
Tables used pandas \parencite{mckinney2010pandas}. Figures used
Matplotlib 3.11.2 \parencite{hunter2007matplotlib}, with seaborn 0.13.2
for regression heatmaps \parencite{waskom2021seaborn}.
The notebook runners recorded nbclient 0.11.0 and nbformat 5.11.1.

Runge row indices are saved in
\nolinkurl{results/runge_part_b_cv/data_split.npz}. The corresponding
folder contains fold and final-fit preprocessing statistics. The
regularisation draws are in
\nolinkurl{results/runge_part_d/configurations.json}.
For AL-PUB, \nolinkurl{results/classification_part_e/reconstructed_manifest.csv}
records dataset membership. That folder also contains
\nolinkurl{class_mapping.json}, \nolinkurl{shared_folds.npz} and
\nolinkurl{sampled_configurations.json}. These files preserve exact
memberships and sampled values alongside the seeded procedures.


The working notebook runs from \nolinkurl{Code/} in the project environment.
Set \nolinkurl{VECLIB_MAXIMUM_THREADS}, \nolinkurl{OMP_NUM_THREADS} and
\nolinkurl{OPENBLAS_NUM_THREADS} to one as shown in the execution guide.

\subsection{Gradient-check pass tolerances}
\label{app:check-tolerances}
For each tested gradient entry, NumPy's \nolinkurl{assert_allclose}
requires $|g-g_{\rm ref}|\leq a+r|g_{\rm ref}|$. Here $a$ is the
absolute tolerance and $r$ is the relative tolerance. The values below
are the pass thresholds in the check scripts, distinct from the
achieved maximum errors reported in Results.
\begin{center}
\begin{tabular}{lrr}
\toprule
Gradient checks & $a$ & $r$ \\
\midrule
Core reuse checks & $10^{-11}$ & $10^{-9}$ \\
Capacity comparison & $10^{-10}$ & $10^{-9}$ \\
Regularisation comparison & $10^{-11}$ & $10^{-9}$ \\
Classification experiments & $2\times10^{-12}$ & $2\times10^{-11}$ \\
\bottomrule
\end{tabular}
\end{center}
The optimiser comparison also has its own saved gradient fixtures. The table specifies the core gradient checks and those for the capacity
comparison, regularisation comparison and classification experiments. Loss values and other behaviours have separate checks and
thresholds in the named scripts. In the penalty checks, bias gradients
must also match the unpenalised path exactly, and zero penalty must
preserve the unregularised training path bitwise.


\subsection{Evidence locations}
Shared code is in \nolinkurl{Code/reused_network.py} and
\nolinkurl{Code/reused_project1.py}. Their comments credit Project 1 and
earlier coursework. Regression settings, fold scores, weights, curves and
predictions are under \nolinkurl{results/runge_part_b_cv},
\nolinkurl{results/runge_part_c} and \nolinkurl{results/runge_part_d}.
Classification settings and folds are under
\nolinkurl{results/classification_part_e}. Per-fit records are in its
\nolinkurl{fits/} folder and final weights, probabilities, confusion matrices
and per-letter scores are in its \nolinkurl{final/} folder.
Historical core check records are in \nolinkurl{results/reuse_checks/checks.json}.
These receipts retain the source hashes from the original runs.

\subsection{Preprocessing illustration and supplementary figures}
\nolinkurl{results/preprocessing_examples/examples.json} records the sampled
training-image paths, publication identifiers and source/prepared-pixel
hashes for the report's preprocessing panel. The panel uses sample seed 2718.
The three supplementary figure exports are in
\nolinkurl{supplementary/figures/}. Matched regularization settings are in
\nolinkurl{supplementary/tables/matched_penalty_settings.csv}, with plotted
pairs in \nolinkurl{results/visual_revision_2026-11-05/matched_penalty_pairs.csv}.
Current classification-table generation uses
\nolinkurl{Code/present_part_e.py}. The notebook contains regression plots.
